\subsection{Translation-conditioned RNA half-life with a decay-factor RBP gate}
\label{sec:q6-translation-conditioned-rna-half-life-rbp-bridge-cross-layer}
\origin{ai}

\paragraph{Finite object.}
The record concerns a finite matched panel of coding-transcript rows.  Each row must carry a synonymous residual coordinate label after quotienting by amino-acid sequence, a same-row translation readout, an RNA half-life readout, measured local occupancy for a prespecified decay-factor RNA-binding protein at the editable RNA window, and matched controls for transcription rate, transcript abundance, splice or isoform state, local RNA structure, motif density, GC and dinucleotide composition, codon-pair context, amino-acid composition, and batch.  The intended contrast is not a free comparison between codons or genes.  It is a controlled comparison between high-coordinate and low-coordinate residual labels, with the half-life statistic interpreted only after the translation readout has been attached to the same row.

\paragraph{Statistic required for a positive bridge.}
Let $i$ range over the matched transcript or CDS rows that survive all joins.  The admissible statistic has the schematic form
$$
\begin{aligned}
\Delta_{HL\mid T}
  &= E[HL_i \mid Q_i=\mathrm{high}, T_i, C_i]
     - E[HL_i \mid Q_i=\mathrm{low}, T_i, C_i],\\
\Delta_{Occ}
  &= E[Occ_i \mid Q_i=\mathrm{high}, C_i]
     - E[Occ_i \mid Q_i=\mathrm{low}, C_i],
\end{aligned}
$$
where $HL_i$ is the RNA half-life readout, $T_i$ is the same-row translation readout, $Occ_i$ is measured local occupancy of the nominated decay-factor RBP, and $C_i$ denotes the declared control vector.  A positive cross-layer record requires more than $\Delta_{HL\mid T}\neq 0$.  The same finite panel must show a local occupancy shift for the nominated RBP, loss of the half-life effect after target-RBP depletion or inactivation, return of the effect under wild-type target-RBP rescue, and failure of rescue under a binding-defective target-RBP or a non-target control.  Without this perturbation-and-rescue pattern, the result remains an RNA-stability association or a coordinate-conditioned descriptor.

\paragraph{Available finite contacts.}
The present payload supplies boundary contacts but not the complete bridge panel.  The RNA half-life contact is organism-scoped and gene- or transcript-keyed; it reports curated half-life estimates or normalized consensus half-life readouts, but it does not include same-window RBP occupancy or perturbation/rescue rows.  The mRNA abundance contact can separate transcription-layer signal from translation-layer signal by supplying matched abundance controls for translation efficiency.  The ribosome-profiling contact supplies per-CDS ribosome occupancy or translation-efficiency estimates, with mRNA abundance used to form $TE=RPF/mRNA$.  These contacts are necessary for the proposed row construction, but they are not sufficient for the decay-factor RBP bridge.

\paragraph{Calibration from neighboring finite records.}
The supporting ledger shows that the project can handle finite cross-layer panels when the required readouts are present.  The RBP-binding calibration contains $900$ total peaks and reports $3$ cross-RBP concordant cases.  The mTOR-translation calibration contains $6494$ rows from the GSE36892 readout source.  The m6Am methylation calibration contains $31886$ rows, $19406$ distinct entries, $20942$ rows in the variable subset, $3230$ low-tail rows below $0.8$, cross-cell concordance marked true, and a saturated-at-least-$0.95$ fraction of $0.343223$.  Its finite statistic reports $p=0.000833$, full-set partial correlation $0.268094$, low-tail partial correlation $0.098666$, variable-subset partial correlation $0.168335$, GC-matched null means $0.003666$, $0.001263$, and $0.002173$ for the full, low-tail, and variable subsets, and positive-control quantities including plus-two-G mean $0.901842$ on $14263$ rows, plus-two-A mean $0.837479$ on $6902$ rows, and SSCA-GC level $0.957884$ on $707$ rows.  Other finite records in the same ledger include $8854$ human and $3990$ yeast rows for the poly(A)-tail comparison, $20000$ methylation probes across $656$ individuals with baseline $R^2=0.0129735$, $657$ N-terminal acetylation rows with AUROC $0.942081$ and $426$ NatA rows, $6288$ protein-acetylation rows with $5921$ distinct stoichiometry entries, $52963$ recombination hotspots, and $2040$ pseudouridylation rows.  These numbers show finite-row discipline, but none of them is the required Q6 RNA-half-life/RBP rescue panel.

\paragraph{Null model.}
The local null is exchangeability of high-coordinate and low-coordinate labels with respect to RNA half-life after conditioning on amino-acid sequence, transcription rate, local RNA structure, motif density, GC and dinucleotide composition, codon-pair context, splice or isoform state, transcript abundance, batch, and same-row translation readout.  A half-life difference without the nominated RBP occupancy shift and the target-RBP loss/rescue pattern does not cross the proposed interpretive boundary.  A model based only on RNA structure, motif density, transcript abundance, isoform state, GC content, dinucleotide content, codon-pair context, or batch is therefore a competitor or control, not support for the RBP bridge.

\paragraph{Local boundary.}
This chapter licenses no translation realization from RNA half-life alone.  It also licenses no protein output, folding, physical admissibility, biological function, causal decay program, phenotype, adaptation, or universal sequence law.  Those claims require their own layer-matched contacts and their own finite statistics.  In the present state, the record is best read as a precise specification of what a positive translation-conditioned RNA-half-life/RBP bridge would have to contain, together with finite neighboring evidence that analogous row-level audits are feasible.  The decisive rows for measured local decay-factor RBP occupancy, target-RBP loss, wild-type rescue, and binding-defective or non-target rescue are not yet attached.
